\documentclass[12pt,reqno,twoside]{amsbook}
%%%%%%%%%%%%%%%%%%%%%%%%%%%%%%%%%%%%%%%%%%%%%%%%%%%%%%%%%%%%
\usepackage{eurosym}
\usepackage{amsmath}
\usepackage{amssymb}
\usepackage{amsfonts}
\usepackage[onehalfspacing]{setspace}
\usepackage{chngcntr}
\usepackage{graphicx}
\usepackage[a4paper, margin=2.5cm]{geometry}
\usepackage[english]{babel}
\usepackage{fancyhdr}
\usepackage{titlesec}
\usepackage{enumitem}
\usepackage{etoolbox}
\usepackage{booktabs}   % nicer tables
\usepackage{hyperref}   % clickable references
\usepackage{csquotes}

\makeatletter
    \def\section{\@startsection{section}{1}%
      \z@{.5\linespacing\@plus.7\linespacing}{.25\linespacing}%
      {\normalfont\bfseries\flushleft}}
    \def\subsection{\@startsection{subsection}{2}%
      \z@{.5\linespacing\@plus.7\linespacing}{.25\linespacing}%
      {\normalfont\bfseries\flushleft}}
\makeatother

\setcounter{MaxMatrixCols}{10}
\numberwithin{section}{chapter}
\fancyhead{}
\fancyfoot{}
\pagestyle{fancy}
\fancyfoot[LE,RO]{\thepage}
\renewcommand{\headrulewidth}{0pt}
\renewcommand{\footrulewidth}{0pt}

\makeatletter
\def\ps@plain{\ps@empty
 \def\@evenfoot{\normalfont\scriptsize\rlap{\thepage}\hfil}%
 \def\@oddfoot{\normalfont\scriptsize\hfil\llap{\thepage}}%
}
\makeatother

%%%%%%%%%%%%%%%%%%%%%%%%%%%%%%%%%%%%%%%%%%%%%%%%%%%%%%%%%%%%
\begin{document}
\frontmatter
\thispagestyle{empty}

% --- Title page ---
\begin{center}
  \vspace*{2cm}
  {\LARGE\bfseries Sentiment Analysis of Movie Reviews\\[0.5em]
   Using Multiple Classification Approaches}\\[1.5em]
  {\large Luiza Coelho \quad Pedro Louro \quad Tiago Vieira}\\[0.5em]
  {\normalsize TMCD 2025/2026 --- 2\textsuperscript{nd} Semester}\\
  {\normalsize Dep.\ de Ciências e Tecnologias da Informação, ISCTE-IUL}\\[1em]
  {\normalsize \today}
\end{center}

\setcounter{page}{1}
\pagenumbering{roman}

% --- Abstract PT ---
\chapter*{Resumo}
\noindent
% TODO: escrever resumo em português (métodos, dados, resultados, conclusão)
% Último parágrafo: contribuição de cada elemento (Luiza: X\%, Pedro: Y\%, Tiago: Z\%)

% --- Abstract EN ---
\chapter*{Abstract}
\noindent
% TODO: English abstract — dataset, methods, key results, conclusions

\tableofcontents

%%%%%%%%%%%%%%%%%%%%%%%%%%%%%%%%%%%%%%%%%%%%%%%%%%%%%%%%%%%%
\mainmatter
\setcounter{page}{1}
\pagenumbering{arabic}

%%%%%%%%%%%%%%%%%%%%%%%%%%%%%%%%%%%%%%%%%%%%%%%%%%%%%%%%%%%%
\chapter{Introduction}
\noindent
% TODO: motivação, descrição da tarefa, visão geral das abordagens, estrutura do relatório

%%%%%%%%%%%%%%%%%%%%%%%%%%%%%%%%%%%%%%%%%%%%%%%%%%%%%%%%%%%%
\chapter{Related Work}
\noindent
% TODO: ≥6 artigos científicos (2×n, n=3)
% Para cada artigo: dataset, abordagem, métricas, comparação com o nosso trabalho

\section{Lexicon and Rule-based Approaches}
\noindent

\section{Machine Learning Approaches}
\noindent

\section{Transformer-based Models}
\noindent

\section{Instruction-based Models}
\noindent

%%%%%%%%%%%%%%%%%%%%%%%%%%%%%%%%%%%%%%%%%%%%%%%%%%%%%%%%%%%%
\chapter{Data}
\noindent
The IMDB reviews dataset~\cite{maas2011learning} consists of movie reviews extracted from
the Internet Movie Database (IMDB), annotated with binary sentiment labels (\textit{pos}/\textit{neg}).

\section{Dataset Characteristics}
\noindent
The dataset is pre-split into train and test sets. Table~\ref{tab:data} summarises the main statistics.

\begin{table}[ht]
\centering
\caption{IMDB reviews dataset statistics.}
\label{tab:data}
\begin{tabular}{lrrrrr}
\toprule
Split & Total & Positive & Negative & Avg.\ tokens & Max tokens \\
\midrule
Train & 21{,}754 & % TODO & % TODO & % TODO & % TODO \\
Test  & 21{,}996 & % TODO & % TODO & % TODO & % TODO \\
\bottomrule
\end{tabular}
\end{table}

% TODO: include fig_class_distribution.png and fig_token_distribution.png

\section{Pre-processing}
\noindent
% TODO: describe preprocessing choices (lowercase, stopword removal, lemmatization, negation)

%%%%%%%%%%%%%%%%%%%%%%%%%%%%%%%%%%%%%%%%%%%%%%%%%%%%%%%%%%%%
\chapter{Methods and Experiments}

\section{Lexicon and Rule-based Models (Task 2.1.1)}
\noindent
% TODO: describe TextBlob, VADER, Stanza and their configuration

\section{Pre-trained Transformer Baseline (Task 2.1.2)}
\noindent
% TODO: describe distilbert-base-uncased-finetuned-sst-2-english and inference setup

\section{NRC Lexicon Classifier (Task 2.2)}
\noindent
% TODO: describe NRC lexicon, scoring approach, negation handling

\section{Classical Machine Learning (Task 2.3)}
\noindent
% TODO: preprocessing pipeline, feature representations, models, hyperparameter search

\section{Fine-tuned Transformer (Task 2.3)}
\noindent
% TODO: distilbert-base-uncased, training setup, computational resources

\section{Instruction-based LLMs (Task 2.4)}
\noindent
% TODO: model, 3 prompts, evaluation subset, prompts in appendix

%%%%%%%%%%%%%%%%%%%%%%%%%%%%%%%%%%%%%%%%%%%%%%%%%%%%%%%%%%%%
\chapter{Results}
\noindent
Table~\ref{tab:results} presents the performance of all approaches on the IMDB test set.

% Paste or \input the LaTeX table from results/table_all_results.tex
% \begin{table}
\caption{Results on IMDB test set}
\label{tab:results}
\begin{tabular}{llrrrr}
\toprule
task & approach & accuracy & precision & recall & f1 \\
\midrule
2.1.1 & Stanza & 0.8335 & 0.9333 & 0.7260 & 0.8167 \\
2.1.1 & VADER & 0.7015 & 0.6604 & 0.8562 & 0.7456 \\
2.1.1 & TextBlob & 0.7000 & 0.6399 & 0.9442 & 0.7628 \\
2.1.2 & DistilBERT (pre-trained, SST-2) & 0.9000 & 0.9228 & 0.8777 & 0.8997 \\
2.2 & NRC Lexicon + Negation & 0.6550 & 0.6254 & 0.8102 & 0.7059 \\
2.2 & NRC Lexicon + Negation & 0.6550 & 0.6254 & 0.8102 & 0.7059 \\
2.2 & NRC Lexicon + Negation & 0.6550 & 0.6254 & 0.8102 & 0.7059 \\
2.2 & NRC Lexicon + Negation & 0.6550 & 0.6254 & 0.8102 & 0.7059 \\
2.2 & NRC Lexicon + Negation & 0.6550 & 0.6254 & 0.8102 & 0.7059 \\
2.2 & NRC Lexicon + Negation & 0.6550 & 0.6254 & 0.8102 & 0.7059 \\
2.2 & NRC Lexicon + Negation & 0.6550 & 0.6254 & 0.8102 & 0.7059 \\
2.2 & NRC Lexicon & 0.6460 & 0.6160 & 0.8160 & 0.7020 \\
2.2 & NRC Lexicon & 0.6460 & 0.6160 & 0.8160 & 0.7020 \\
2.2 & NRC Lexicon & 0.6460 & 0.6160 & 0.8160 & 0.7020 \\
2.2 & NRC Lexicon & 0.6460 & 0.6160 & 0.8160 & 0.7020 \\
2.2 & NRC Lexicon & 0.6460 & 0.6160 & 0.8160 & 0.7020 \\
2.2 & NRC Lexicon & 0.6460 & 0.6160 & 0.8160 & 0.7020 \\
2.2 & NRC Lexicon & 0.6460 & 0.6160 & 0.8160 & 0.7020 \\
2.3 & DistilBERT (fine-tuned) & 0.9475 & 0.9508 & 0.9462 & 0.9485 \\
2.3 & SVM (TF-IDF bi + neg) & 0.9045 & 0.9094 & 0.9031 & 0.9062 \\
2.3 & SVM (TF-IDF bi + neg, C=0.1) & 0.9035 & 0.9060 & 0.9051 & 0.9055 \\
2.3 & LR (TF-IDF bi) & 0.9000 & 0.9014 & 0.9031 & 0.9022 \\
2.3 & SVM (TF-IDF uni) & 0.9000 & 0.9014 & 0.9031 & 0.9022 \\
2.3 & LR (TF-IDF uni) & 0.8925 & 0.8929 & 0.8973 & 0.8951 \\
2.3 & LR (BoW) & 0.8825 & 0.8885 & 0.8806 & 0.8845 \\
2.3 & LR (BoW) & 0.8815 & 0.8875 & 0.8796 & 0.8835 \\
2.3 & NB (BoW) & 0.8580 & 0.8758 & 0.8415 & 0.8583 \\
2.3 & SVM (BoW) & 0.8540 & 0.8614 & 0.8513 & 0.8563 \\
2.4 & Claude Haiku (claude-haiku-4-5-20251001) — domain prompt & 0.9550 & 0.9417 & 0.9700 & 0.9557 \\
2.4 & Claude Haiku (claude-haiku-4-5-20251001) — generic prompt & 0.9500 & 0.9412 & 0.9600 & 0.9505 \\
2.4 & Claude Haiku (claude-haiku-4-5-20251001) — few_shot prompt & 0.9500 & 0.9412 & 0.9600 & 0.9505 \\
\bottomrule
\end{tabular}
\end{table}


% TODO: analysis and discussion of results
% TODO: include fig_all_results.png

%%%%%%%%%%%%%%%%%%%%%%%%%%%%%%%%%%%%%%%%%%%%%%%%%%%%%%%%%%%%
\chapter{Conclusions}
\noindent
% TODO: best approach, limitations, future work

%%%%%%%%%%%%%%%%%%%%%%%%%%%%%%%%%%%%%%%%%%%%%%%%%%%%%%%%%%%%
\renewcommand{\bibname}{References}
\def\bibindent{0.7cm}
\begin{thebibliography}{99\kern\bibindent}
\makeatletter
\let\old@biblabel\@biblabel
\def\@biblabel#1{\old@biblabel{#1}\kern\bibindent}
\let\old@bibitem\bibitem
\def\bibitem#1{\old@bibitem{#1}\leavevmode\kern-\bibindent}
\makeatother
\makeatletter
\renewcommand\@biblabel[1]{}
\makeatother

\bibitem{maas2011learning}
  A.~L. Maas, R.~E. Daly, P.~T. Pham, D.~Huang, A.~Y. Ng, and C.~Potts (2011),
  \textquotedblleft Learning Word Vectors for Sentiment Analysis\textquotedblright,
  in \textit{Proceedings of ACL}, pp.\ 142--150.

% TODO: add ≥5 more references

\end{thebibliography}

%%%%%%%%%%%%%%%%%%%%%%%%%%%%%%%%%%%%%%%%%%%%%%%%%%%%%%%%%%%%
\appendix
\chapter*{Appendix A — LLM Prompts}
\noindent

\paragraph{Prompt 1 — Generic}
\begin{quotation}
\noindent Classify the sentiment of the following text as ``positive'' or ``negative''.
Reply with only one word.\\
Text: \{review\}
\end{quotation}

\paragraph{Prompt 2 — Domain-aware}
\begin{quotation}
\noindent You are analyzing movie reviews. Classify the sentiment of the following
movie review as ``positive'' or ``negative''. Reply with only one word.\\
Review: \{review\}
\end{quotation}

\paragraph{Prompt 3 — Few-shot}
\begin{quotation}
\noindent Classify movie review sentiment as ``positive'' or ``negative''.\\[0.3em]
Examples:\\
{[POS]} \{example\_pos\_1\}\\
{[NEG]} \{example\_neg\_1\}\\
{[POS]} \{example\_pos\_2\}\\
{[NEG]} \{example\_neg\_2\}\\
{[POS]} \{example\_pos\_3\}\\
{[NEG]} \{example\_neg\_3\}\\[0.3em]
Now classify: \{review\}\\
Reply with only one word: positive or negative.
\end{quotation}

\end{document}
